\documentclass[12pt]{article}
\usepackage[T1]{fontenc}
\usepackage[utf8]{inputenc}
\usepackage{lmodern}
\usepackage{textcomp}
\usepackage{enumitem}
\usepackage{geometry}
\usepackage{graphicx}
\usepackage{amsmath, amssymb}
\geometry{margin=1in}
\begin{document}
\begin{center}
Astronomy - Graduate Level Assessment 21 \
Total Marks: 40 \
Answer all questions.
\end{center}

\section*{Multiple Choice (10 marks)}
\begin{enumerate}[label=\arabic*.]
\item You cool a blackbody to half its original temperature. How does its spectrum change?
\begin{enumerate}[label=(\alph*)]
    \item Power emitted is 1/16 times as high; peak emission wavelength is 1/2 as long.
    \item Power emitted is 1/4 times as high; peak emission wavelength is 2 times longer.
    \item Power emitted is 1/4 times as high; peak emission wavelength is 1/2 as long.
    \item Power emitted is 1/16 times as high; peak emission wavelength is 2 times longer.
\end{enumerate}
\item A sand bag has a mass of 5 kg and weight 50 N on Earth. What is the mass and weight of this sand bag on a planet with half the gravity compared to Earth?
\begin{enumerate}[label=(\alph*)]
    \item Mass 5 kg weight 100 N
    \item Mass 5 kg weight 50 N
    \item Mass 5 kg weight 25 N
    \item Mass 10 kg weight 100 N
\end{enumerate}
\item Our current best observations show that Pluto has
\begin{enumerate}[label=(\alph*)]
    \item one medium sized satellite and two small satellites.
    \item no satellites.
    \item one large satellite and three small satellites.
    \item one large satellite.
\end{enumerate}
\item Mars has an atmosphere that is almost entirely carbon dioxide.Why isn't there a strong greenhouse effect keeping the planet warm?
\begin{enumerate}[label=(\alph*)]
    \item Mars does not have enough internal heat to drive the greenhouse effect
    \item Mars is too far from the sun for the greenhouse effect to work
    \item the greenhouse effect requires an ozone layer which Mars does not have
    \item the atmosphere on Mars is too thin to trap a significant amount of heat
\end{enumerate}
\item The famous Drake equation attempts to answer the following question:
\begin{enumerate}[label=(\alph*)]
    \item Will the Sun become a black hole?
    \item Is the universe infinitely large?
    \item How old is the visible universe?
    \item Are we alone in the universe?
\end{enumerate}
\end{enumerate}

\section*{True / False (10 marks)}
\textit{Answer True or False}
\begin{enumerate}[label=\arabic*.]
\setcounter{enumi}{5}
\item True or False: A black body is an idealized physical object that reflects all electromagnetic radiation without absorption.
\item True or False: The Phoenix Mars Lander successfully landed on Mars on May 25, 2008, to investigate ice-rich soil in the polar region.
\item True or False: A scientific theory can be conclusively proven correct after a single successful experiment under controlled conditions.
\item True or False: The Large Magellanic Cloud is the closest planetary nebula to Earth.
\item True or False: Venus does not experience seasons because its rotation axis is nearly perpendicular to the plane of the Solar System.
\end{enumerate}

\section*{Long Form Response (20 marks)}
\textit{Answer each question with a clear and complete explanation.}
\begin{enumerate}[label=\arabic*.]
\setcounter{enumi}{10}
\item Discuss the differences between sidereal and solar days, and critically analyze the statement that the Moon's circuit in our sky corresponds to one solar day.
\item Discuss the processes that led to the differentiation in material composition between the inner and outer planets in our solar system, focusing on the role of temperature in the nebula.
\end{enumerate}

\vfill
\noindent\textit{End of Paper}
\end{document}